\documentclass{article}

% El formato oficial del concurso es la plantilla NeurIPS. Si neurips_2026.sty
% está en este directorio se usa; si no, se cae a una geometría equivalente para
% que el documento compile en cualquier máquina.
\IfFileExists{neurips_2026.sty}{\usepackage[final]{neurips_2026}}{%
  \usepackage[margin=1in]{geometry}
  \usepackage{times}
}

\usepackage[utf8]{inputenc}
\usepackage[T1]{fontenc}
\usepackage{hyperref}
\usepackage{url}
\usepackage{booktabs}
\usepackage{amsmath}
\usepackage{amssymb}
\usepackage{microtype}

\title{Shrinkage improves predictive AI evaluation\\
under benchmark holdout:\\
a reproducible study and public-data audit}

\author{%
  Francisco Angulo de Lafuente\\
  Independent researcher, Madrid, Spain\\
  \texttt{lareliquia.angulo@gmail.com}\\
}

\begin{document}
\maketitle

\begin{abstract}
We study predictive evaluation of AI systems under the benchmark-level holdout
defined by the Predictive AI Evaluation Competition, where the items to be
predicted come from benchmarks absent from training. From the public
\texttt{measurement-db} release we build a local simulation of the competition's
Brier ALC, validate it against hand-constructed cases, and evaluate under
leave-one-benchmark-out over 101 held-out public benchmarks that were not
inspected during development. Our submission is the posterior mean of a
Beta--Binomial model with a symmetric prior and a single shrinkage strength,
taken as the argmin of a disjoint 50-benchmark development split. It reduces
Brier ALC from $0.2454$ to $0.2105$ against the official empirical-mean
reference: an absolute reduction of $0.0349$ and a relative reduction of
$14.2\%$. On a run with a randomised acquisition order, resampling benchmarks with their
pairs kept together gives $-0.0310$, 95\% CI $[-0.0373,-0.0249]$. The gain improves 86 of 101 benchmarks,
survives dropping the five benchmarks where it is largest, holds between
$13.1\%$ and $14.5\%$ across every sampling choice we varied, and falls only to
$-0.0323$ when benchmarks are reweighted by their eligible pair counts. All of
these figures come from our own approximation of the evaluation, not from the
official scorer, and the reference they beat is deliberately simple.
We then report seven extensions that produced no detectable improvement, including
two direct attempts at the largest remaining target: an oracle restricted to one
value per pair scores $0.1789$, and $96.7\%$ of that floor is variation
\emph{between items inside a pair}, yet models of per-item deviation fitted
out-of-benchmark --- with hand-built features and with hashed sparse text ---
made results strictly worse at every weight we tried, even when started from the
true pair rate.
Finally we audit the release and report three quantified findings. All code, the
harness and its ten tests are released.
\end{abstract}

\section{Task and scoring}

A \emph{subject} is an AI system fixed by model, harness, harness version and
reasoning effort. An \emph{item} is a task. A \emph{response} is one scored run,
$1$ if correct. The predictor receives the subject attributes, the item
attributes (content, features, interactors, an anonymous \texttt{benchmark\_id})
and whatever labels the current budget has revealed, and returns
$\hat p \in [0,1]$ without running the subject.

Two properties of the scoring rule drive what follows.

\textbf{The unit of scoring is the subject--benchmark pair.}
Brier ALC is computed per pair and the pairs are averaged with equal weight,
regardless of how many responses each contributes:
\begin{equation}
\mathrm{ALC} = 0.1\,B_0 + 0.2\,B_1 + 0.2\,B_3 + 0.2\,B_7 + 0.2\,B_{15} + 0.1\,B_{31},
\label{eq:alc}
\end{equation}
where $B_n$ is the Brier score at a budget of $n$ revealed labels for that pair.
Section~\ref{sec:negativos} shows that matching a prior's averaging to
Equation~\eqref{eq:alc} mattered more, in our experiments, than any modelling
refinement we tried.

\textbf{The holdout is at benchmark level.} Test items are drawn from benchmarks
assigned to a private pool, so benchmark-specific parameters cannot be
transferred. Whether \emph{item} properties transfer is an empirical question;
we test it twice in Section~\ref{sec:negativos}, for absolute difficulty and for
within-pair relative difficulty. We validate by holding out whole benchmarks and
never by splitting responses at random: $68.8$M observations are not $68.8$M
independent examples of transfer.

\section{Data}

The public release contains 249 curated benchmarks. Of the 281 catalogue records
we read, 278 are marked \texttt{release == public} (three are withheld for source
reproducibility) and 165 of those are binary, contributing $68{,}840{,}488$
recorded responses over $14{,}733$ subject--benchmark rows, of which $13{,}289$
have at least 80 responses. We use the \texttt{responses}, \texttt{items},
\texttt{subjects} and \texttt{benchmarks} tables; \texttt{traces} and
\texttt{assets} ($244$\,GiB together) are never downloaded because the predictor
is never shown them. Everything in this paper concerns those 165 binary
benchmarks, not the release as a whole.

\paragraph{Local evaluation set.} We follow the published construction: pairs with
at least 80 distinct items, items split 50/50 into an acquisition pool and an
evaluation pool with a fixed per-pair seed, acquisition items never scored,
evaluation outcomes never revealed, nested label sets at budgets
$\{0,1,3,7,15,31\}$, and one \texttt{predict} call per distinct evaluation item
whose value is reused by every recorded response of that item. Two choices are
ours and not part of the official procedure: we sample at most six pairs per
benchmark and cap each pair at 150 distinct evaluation items.
Section~\ref{sec:sens} measures how much they matter. The result is 769 pairs
over 151 benchmarks and $4.78$M scored responses.

\paragraph{Development and headline splits.} From a fixed seed we assigned 50
benchmarks to development and 101 to the headline set (496 pairs). The shrinkage
strength and the choice of submitted method were fixed on development.
Section~\ref{sec:limites} states exactly which analyses touched the headline set
and when.

\section{Harness validation}

We have neither the official scorer nor the private data, so we do not claim
numerical agreement with the official score. We claim that our implementation
satisfies the published definition, and we test it. Ten checks ship with the code
(\texttt{src/test\_metrica.py}), covering the weights, the constant-$0.5$ and
perfect-predictor limits, pool disjointness, nesting, the 80-item threshold,
equal pair weighting, reuse of one prediction across repeated responses, and the
hand-built case where unsmoothed averaging scores $0.925$. All ten pass.

All partitions are derived with \texttt{blake2b}, not Python's built-in
\texttt{hash}, which is randomised per process. An earlier version of this work
used the latter; an independent re-derivation caught it, and every number below
was regenerated afterwards. We mention it because it is the kind of defect that
silently makes a reproducibility claim false.

\section{Method}

Our submission is the posterior mean of a Beta--Binomial model for the pair:
\begin{equation}
\hat p \;=\; \frac{\alpha + k}{\alpha + \beta + n},
\qquad \alpha = P_0 S, \quad \beta = (1-P_0) S,
\label{eq:bb}
\end{equation}
with $k$ successes among the $n$ labels revealed for that subject and that
\texttt{benchmark\_id}, $P_0 = 0.5$ and $S = 3$. $S=3$ is the argmin on
development inside a region that is flat over $S\in[3,5]$
($0.21708$, $0.21740$, $0.21809$). The predictor keeps no state, fits nothing at
run time, and needs no network, no files and no payload. If
\texttt{benchmark\_id} is ever absent it returns the prior rather than pooling
labels across benchmarks.

The official reference instead returns the raw mean of the revealed labels, and
$0.5$ when none exist. At $n=1$ that is $0$ or $1$, so a single wrong guess costs
the maximum Brier loss of $1$. This is not hypothetical: on the headline set the
reference scores $B_1 = 0.3650$, worse than ignoring the labels entirely, and
Table~\ref{tab:main} shows that almost the whole gap between the two methods
opens at $B_1$ and $B_3$.

\section{Results}

\begin{table}[t]
\centering
\caption{Headline set: 101 held-out benchmarks, 496 pairs, default acquisition
order (\texttt{src/ejecuta.py final}). Lower is better. The three rows nearest
the submission are within noise of it; none was chosen after seeing this table.}
\label{tab:main}
\small
\begin{tabular}{lrrrrrrr}
\toprule
Predictor & ALC & $B_0$ & $B_1$ & $B_3$ & $B_7$ & $B_{15}$ & $B_{31}$ \\
\midrule
Empirical-Bayes prior, subject-conditioned & 0.2099 & 0.2471 & 0.2311 & 0.2119 & 0.1991 & 0.1905 & 0.1865 \\
Hierarchical prior, pair-averaged          & 0.2102 & 0.2484 & 0.2313 & 0.2127 & 0.1994 & 0.1905 & 0.1863 \\
\textbf{Beta--Binomial, $S{=}3$ (submitted)} & \textbf{0.2105} & 0.2500 & 0.2319 & 0.2128 & 0.1992 & 0.1905 & 0.1863 \\
Empirical-Bayes prior                      & 0.2109 & 0.2512 & 0.2332 & 0.2125 & 0.1992 & 0.1906 & 0.1864 \\
Ablation: subject effect only              & 0.2137 & 0.2650 & 0.2368 & 0.2156 & 0.1996 & 0.1906 & 0.1863 \\
Ablation: global mean only                 & 0.2175 & 0.2770 & 0.2450 & 0.2188 & 0.2007 & 0.1912 & 0.1864 \\
Empirical mean (\emph{official reference}) & 0.2454 & 0.2500 & \textbf{0.3650} & 0.2434 & 0.2076 & 0.1925 & 0.1868 \\
Constant $0.5$                             & 0.2500 & 0.2500 & 0.2500 & 0.2500 & 0.2500 & 0.2500 & 0.2500 \\
\midrule
\emph{Oracle, one value per pair}          & \emph{0.1789} & & & & & & \\
\bottomrule
\end{tabular}
\end{table}

Table~\ref{tab:main} gives the headline result: Brier ALC falls from $0.2454$ to
$0.2105$, an absolute reduction of $0.0349$ and a relative reduction of $14.2\%$.

\paragraph{Uncertainty.} The 496 pairs are not 496 independent units for an
inference about new benchmarks, because pairs cluster inside their benchmark. We
therefore resample \emph{benchmarks}, keeping their pairs together. On a run with
a randomised acquisition order the difference is $-0.0310$ with a 95\% interval
of $[-0.0373,-0.0249]$ and no replicate of the opposite sign. Averaging per
benchmark rather than per pair gives $0.2088$ against $0.2393$ on that run; the
two units agree closely here, but we report them separately because they are not
interchangeable.

\paragraph{Where the gain lives.} Back on the default-order run, the method
improves 86 of 101 benchmarks and
worsens 15. The median per-benchmark difference is $-0.0405$, the 10th and 90th
percentiles are $-0.0701$ and $+0.0127$, the best case is $-0.1273$ and the worst
$+0.0463$. Dropping the five benchmarks where it gains most moves the mean
difference only from $-0.0347$ to $-0.0317$, so the result is not carried by a
handful of benchmarks.

\paragraph{Stability.} Across five random acquisition draws our score ranges over
$[0.2096,0.2110]$ (s.d.\ $0.0005$) and the reference over $[0.2406,0.2447]$
(s.d.\ $0.0014$).

\paragraph{How much is left.} An oracle that knows each pair's accuracy but must
emit one value per pair --- the best any pair-constant predictor can do, ours and
the reference included --- scores $0.1789$ on the same pairs. The reference leaves
$0.0665$ of headroom and we capture $0.0349$, or $52.5\%$. Decomposing that floor,
$0.0059$ ($3.3\%$) is variation between repeated runs of the same item and
$0.1730$ ($96.7\%$) is variation \emph{between items inside a pair}. The rest of
the problem is therefore about which items in an unseen benchmark are hard, and
Section~\ref{sec:negativos} reports our failed attempt on exactly that.

\paragraph{A claim we decline to make.} It would be convenient to say the official
reference is no better than doing nothing. It is better: mean ALC $0.2454$ against
$0.2500$. But the benchmark-cluster interval for that difference,
$[-0.0064,+0.0252]$ on the randomised-order run, includes zero; with the default
order the gap is $+0.0046$. Our $14.2\%$ relative reduction should be read
alongside the absolute difference of $0.0349$, and against a reference that is
deliberately simple rather than competitive.

\section{Sensitivity to our local approximations}
\label{sec:sens}

Varying the cap on evaluation items per pair over $\{25,50,100,150\}$ gives
relative reductions of $14.5\%$, $14.3\%$, $14.3\%$ and $14.2\%$. Varying the
number of pairs sampled per benchmark over $\{1,2,4,6\}$ gives $13.1\%$, $14.3\%$,
$14.2\%$ and $14.2\%$.

A third approximation is less innocent. Sampling a fixed number of pairs per
benchmark gives every benchmark equal weight, whereas the competition averages
pairs. Among the headline benchmarks there are $11{,}657$ eligible pairs and the
five largest hold $77.9\%$ of them (\texttt{alignment\_faking} $36.7\%$,
\texttt{sib200} $28.7\%$). Reweighting the per-benchmark differences by eligible
pair counts moves the mean difference from $-0.0347$ to $-0.0323$: still clearly
favourable, but we report it because our proxy does under-weight the benchmarks
the official metric would emphasise.

\section{Extensions that produced no detectable improvement}
\label{sec:negativos}

\paragraph{Averaging the prior the wrong way cost more than any refinement.}
Our first hierarchical prior used the response-weighted global accuracy, $0.350$,
which is dominated by a handful of enormous benchmarks. The quantity
Equation~\eqref{eq:alc} averages is the per-pair accuracy, $0.538$ among the rows
the competition evaluates. Using the
former raised ALC to $0.2137$ and $B_0$ to $0.2650$, against $0.2105$ and
$0.2500$ in the same run. Among the eligible rows the two averages are $0.350$ and $0.538$; the figure of
$0.549$ quoted above is over all rows. For this objective, averaging the prior the way the
score is averaged was the single most consequential choice we made; we do not
offer it as a general rule for conditional priors.

\paragraph{A subject-identity prior does not transfer reliably.} Conditioning the
prior on the subject's accuracy in other benchmarks improves $B_0$ from $0.2500$
to $0.2471$ --- a real if small signal --- but the ALC difference against the
plain symmetric prior is $-0.0005$ with a 95\% interval of $[-0.0024,+0.0012]$ and
$p=0.64$ on the randomised-order run. With the default acquisition order it wins
in 46 of 101 benchmarks, sign-flip $p=0.43$. Among eligible rows only 16 subjects appear in
five or more benchmarks, and a subject's accuracy varies by $0.130$ across
benchmarks against a between-pair spread of $0.273$
(\texttt{src/estadisticas\_datos.py}).

\paragraph{A non-parametric prior over the bimodal distribution does not help.}
Per-pair accuracy is strongly bimodal. Among the $13{,}289$ rows with at least
80 responses --- the only ones the competition evaluates --- the modal bin is
$[0,0.05)$, holding $8.3\%$ of them, and a broad second mode covers the upper
half; the macro mean is $0.538$ with a spread of $0.273$. A single Beta cannot
represent that shape (\texttt{src/estadisticas\_datos.py}). A
non-parametric prior on a 101-point grid, fitted out-of-benchmark, scores
$+0.0001$ against Equation~\eqref{eq:bb} ($p=0.81$ on the randomised-order run,
$-0.00002$ and $p=0.55$, 47 of 101 benchmarks, with the default order), and this held across filters
on minimum pair size, recentring, and uniform mixing weights from $0.05$ to
$0.3$ (\texttt{src/barrido\_prior.py}). The experiment does not establish whether
the limitation lies in estimating the prior, in transfer across benchmarks, or in
the information the revealed labels carry.

\paragraph{Per-budget shrinkage tuning buys nothing.} Because the predictor is
constant per pair, each budget can be optimised independently. Doing so selects
$S=3$ at every budget with at least one label and reaches $0.21708$ on
development, identical to a single $S=3$. At budget $0$ there are no labels, so
every $S$ gives the same prediction and the sweep is degenerate there.

\paragraph{Cheap item-text features give nothing detectable at budget 0.} We
fitted a regularised logistic model on $1.17$M development observations over 23
columns: twelve item-text features computable inside \texttt{predict} --- length, word count, digit,
uppercase, non-ASCII and punctuation ratios, code and \LaTeX{} markers, option
count, multiple-choice indicator, question marks, line breaks --- plus release
year, harness and provider indicators. Under leave-one-benchmark-out
\emph{within development} it
scores $0.2495$ against $0.2500$ for the constant, winning in 18 of 50 benchmarks
($p = 0.065$), worth $-0.00005$ in ALC. Several coefficients have plausible signs
(multiple-choice easier, code and length harder), but coefficient signs alone do
not establish transferable signal.

\paragraph{Per-item prediction, the largest target, made things worse.} Since
$96.7\%$ of the remaining floor is between-item variation and \texttt{predict}
receives the specific item, a pair-constant predictor leaves that entire share
untouched by construction. The previous experiment asked whether features predict
\emph{absolute} correctness, which conflates pair difficulty with item
difficulty. We therefore asked the sharper question: do they predict an item's
\emph{deviation from its own pair's mean}? We fitted a ridge model on $1.63$M
development observations with the pair mean as an offset, so the coefficients
capture only within-pair relative difficulty, and added its prediction to
Equation~\eqref{eq:bb} with weight $\lambda$. Development ALC was $0.21708$ at
$\lambda=0$, matching Equation~\eqref{eq:bb} exactly, and rose monotonically to
$0.21776$, $0.21896$ and $0.22067$ at $\lambda = 0.5, 1.0, 1.5$. Adding this
signal is strictly harmful at every weight we tried.

\paragraph{Nor does a far richer text representation.} To bound that result
beyond one hand-built feature set, we repeated it with hashed sparse text
features --- $2^{18}$ word uni/bigrams plus $2^{17}$ character 3--4-grams, fitted
with SGD --- and made the test as generous as possible by starting from the
\emph{true} pair mean rather than from our estimate. Over 50 development
benchmarks the oracle pair mean scores $0.1905$; adding the predicted deviation
gives $0.1907$ at $\lambda=0.5$ and $0.1912$ at $\lambda=1.0$, improving only 14
of 50 benchmarks (sign-flip $p=0.003$ \emph{against} the addition). The mean
correlation between predicted and realised deviation is $-0.016$, with a maximum
of $0.196$ on a single benchmark. So even given the pair rate for free, cheap text
features do not identify which items an unseen benchmark's subject will miss.

A route we could not test: the organisers' companion embedding release
\texttt{measurement-db-embed} covers only 3 of the 165 binary benchmarks, which is
too few to measure transfer across benchmarks. A learned representation from a
language model remains the obvious untried candidate.

\paragraph{Why we submit no acquisition function.} The competition allows an
optional \texttt{labeling.py} to choose which outcomes to reveal. We submit none,
and we want to be careful about why. Our predictor cannot distinguish items, so it
has no basis on which to prefer one candidate over another, and we fall back to
the platform's random default. We do \emph{not} claim random is optimal: since
$96.7\%$ of the pair-constant floor is between-item variance, a policy that
selected candidates by content could in principle estimate the pair rate with
less variance than a random draw. That is an untested opportunity we leave open,
and it rests on the same item-level signal the previous experiment failed to find.

Taken together, in this protocol and on this release, the improvement over the
official reference comes from calibrating how far a few revealed labels may move
the estimate, and seven attempts to add information beyond that produced nothing
detectable. This bounds what we extracted with cheap, fully offline methods.
Item embeddings, inference-time model calls and richer text representations
remain untested, and the $0.1730$ of between-item variance remains available to
whoever can reach it.

\section{Audit of the public release}

Definitions were fixed before looking at the data: ``same item'' means an
identical published \texttt{content\_hash}; ``repetition'' means two response rows
with the same subject and item; a ``post-response field'' describes how an answer
is graded rather than what the system was shown. Nothing here uses platform
feedback, and we made no attempt to identify the subjects or benchmarks behind
anonymous identifiers.

\paragraph{Confirmed: items are shared across benchmarks.} Of $2{,}898{,}062$
distinct content hashes over $3{,}770{,}701$ items in the 165 binary benchmarks,
$1{,}980$ ($0.068\%$) appear in more than one benchmark, involving 11 benchmarks.
The largest overlaps are \texttt{adaptivestep}$\leftrightarrow$\texttt{prox}
($1{,}819$ items), \texttt{adaptivestep}$\leftrightarrow$\texttt{felm} and
\texttt{felm}$\leftrightarrow$\texttt{prox} ($215$ each), and
\texttt{swe\_rebench}$\leftrightarrow$\texttt{swebench\_live} ($91$). This is
irrelevant to our submitted predictor, which never reads item content, but not in
general: any method using item text could see the same item in two benchmarks and
count it as independent evidence. The case that would matter most --- whether any
such item also sits in a private test benchmark --- cannot be checked from
outside, and we do not claim it does. We raise it as a question for the
organisers, who can.

\paragraph{Not found: leakage into predictor-visible fields.} We screened the
fields the predictor actually receives for answer-, solution-, label- and
grading-like content across 60 of the 165 benchmarks. Two had keyword matches in
\texttt{item\_features}, and in neither case did the field coincide with the
benchmark's grading criterion. We record this as no finding in a keyword screen
over 60 benchmarks, which is not evidence of absence.

\paragraph{Confirmed: a fifth of repeated observations contain both outcomes.}
Across 87 benchmarks with repeated measurements there are $3{,}154{,}268$
subject--item pairs measured more than once, and $635{,}977$ of them
($20.2\%$; median $8.3\%$ per benchmark) contain both a success and a failure.
Some pairs are repeated up to $10{,}000$ times. This need not indicate a grader
fault --- stochastic sampling by the system explains much of it --- and it does
\emph{not} by itself explain the $0.1789$ pair-constant oracle score: as measured
in Section~\ref{sec:sens}, repetition variance accounts for $0.0059$ of that
figure and between-item variance for $0.1730$. What it does mean is that a
measurable share of response-level Brier loss is irreducible for any predictor,
and that should be stated whenever a Brier number on this data is interpreted.

\section{Limitations}
\label{sec:limites}

Our local score is not the official score. We approximate the evaluation with our
own harness, cap evaluation items at 150 per pair, sample at most six pairs per
benchmark, restrict to binary benchmarks, and weight benchmarks rather than pairs
(Section~\ref{sec:sens} quantifies the last three). We report Brier ALC on public
training benchmarks held out from each other, which is a proxy for, not a
measurement of, performance on the private pool.

On development discipline we want to be precise rather than flattering. The
shrinkage strength is the development argmin, written by
\texttt{src/barrido\_prior.py}, and we did not select it by looking at the
headline set. We did, however, score the headline set twice with two values of
it. An earlier version of this work fixed $S=4$ by hand inside the flat region and
reported $0.2111$ and a $14.0\%$ reduction; we then replaced it with the
development argmin $S=3$, which gives $0.2105$ and $14.2\%$. The change was made
to remove an arbitrary choice, not because the headline improved, and it moves the
result by $0.0006$. We state it because a reader cannot otherwise tell from the
artefacts that two headline runs exist. Separately, the headline set was not
scored exactly once: it was
scored by the final run, then again under five acquisition draws and
$2{,}000$ bootstrap replicates, and the sensitivity study re-scored it under
other sampling choices. The per-budget sweep, the two text experiments and the
sensitivity study were all designed after the headline numbers existed; all are
reported on development or as sensitivity around an already-fixed result, and
none changed the submission, but they are exploratory rather than
pre-registered. The audit statistics and the bimodality figures are computed over
all 165 binary benchmarks, headline ones included, since they describe the data
rather than evaluate a method. An intermediate version of the headline table was
also discussed with a language model used as a critic during development; that
discussion changed the framing and the analyses, not the submitted predictor.

Our negative results are bounded by what we tried, with cheap offline methods, on
this release, under our approximation of the protocol.

\section{Reproducibility}

Everything in this paper is produced by a named script. The data are the public
\texttt{measurement-db} release, fetched with \texttt{src/descarga\_datos.py}, at
verified data commit
\begin{center}\small\texttt{a23e041be8651be4d97decba385779300abeabee}\end{center}
and builder revision
\begin{center}\small\texttt{287495967fc276275c3b466dfede2cc3534e830c}.\end{center} The evaluation set is built by \texttt{src/construye.py} with master seed
$20261024$, and the development/headline partition is drawn from the same seed in
\texttt{src/ejecuta.py}. Every result maps to one script:

\begin{center}\small
\begin{tabular}{ll}
\toprule
Result & Script \\
\midrule
Ten harness checks                       & \texttt{test\_metrica.py} \\
Descriptive statistics                   & \texttt{estadisticas\_datos.py} \\
Submitted shrinkage strength (writes \texttt{hiper.json}) & \texttt{barrido\_prior.py} \\
Per-budget sweep                         & \texttt{barrido\_presupuesto.py} \\
Hierarchical-variant curve               & \texttt{ejecuta.py tune} \\
Table~\ref{tab:main}                     & \texttt{ejecuta.py final} \\
Uncertainty, stability, both averaging units & \texttt{robustez.py} \\
Sensitivity, oracle decomposition, reweighting & \texttt{sensibilidad.py} \\
Audit                                    & \texttt{auditoria.py} \\
Text experiments                         & \texttt{texto\_b0.py}, \texttt{por\_item.py} \\
\bottomrule
\end{tabular}
\end{center}

\noindent The \texttt{tune} mode sweeps the hierarchical variants only and does
not write \texttt{hiper.json}. The submitted strength is reproduced by running
\texttt{barrido\_prior.py}, the sole writer of that file.
Requirements are Python 3.10+, \texttt{pandas},
\texttt{pyarrow}, \texttt{numpy}, \texttt{scikit-learn} and
\texttt{huggingface\_hub}, about $14$\,GiB of disk, and a few hours on one
desktop CPU with no GPU. The submitted \texttt{model.py} is standard library
only, carries no payload and depends on nothing at run time.

\paragraph{Disclosure.} All training and development used only the public
\texttt{measurement-db} release listed in the benchmark inventory; no other data
source was used. The work was carried out with substantial assistance from large
language models, for code, for prose and as adversarial critics of the analysis.
All numerical results were produced by the released scripts, and an independent
re-derivation --- which found and forced the correction of the non-deterministic
partitioning described in Section 3, among other errors --- was run before
submission.

\end{document}
